\begin{recipe}{Spinazierisotto met Parmezaanse kaas}
\kicker[Hoofdgerecht,Vegetarisch,Italiaans]{Hoofdgerecht · vegetarisch · Italiaans}
\index[register]{Spinazierisotto met Parmezaanse kaas}
\index[register]{Spinazie}
\index[register]{Parmezaanse kaas}
\index[register]{Risottorijst}

\meta{35 minuten}

\lettrine{D}{e} spinazie gaat er pas op het eind bij, net voor de risotto klaar
is, zodat hij net slinkt en zijn kleur en frisheid houdt. De boter erdoor
roeren aan het einde, de mantecatura, maakt de risotto extra romig.

\blockrule{Ingrediënten}
\begin{ingredients}
  \ing{200--250 g}{verse spinazie}\index[register]{Spinazie}
  \ing{60 g}{Parmigiano Reggiano}\index[register]{Parmezaanse kaas}
  \ing{320 g}{risottorijst}\index[register]{Risottorijst}
  \ing{1 liter}{warme bouillon}
  \ing{150 ml}{droge witte wijn}
  \ing{1}{kleine ui, gesnipperd}
  \ing{30--40 g}{boter}
  \ingb{scheutje olijfolie, om te fruiten}
  \ingb{extra Parmezaanse kaas, om te serveren}
\end{ingredients}

\blockrule{Bereiding}
\tip{De mantecatura, de boter (en kaas) aan het eind door de risotto
kloppen met de hapjespan van het vuur, maakt het verschil tussen een gewone en
een romige, glanzende risotto.}
\begin{steps}
  \item Snipper de ui en breng de bouillon aan de kook. Zet daarna het vuur
        laag zodat de bouillon warm blijft.
  \item Fruit de ui in een scheutje olijfolie in een hapjespan met dikke
        bodem tot hij glazig is.
  \item Voeg de risottorijst toe en rooster hem al roerend kort mee.
        Blus af met de witte wijn en roer tot die is opgenomen.
  \item Voeg de bouillon in kleine scheutjes toe, telkens wachtend tot die
        is opgenomen voor je de volgende schep toevoegt, tot de rijst bijna
        gaar is, ongeveer 18--20 minuten.
  \item Roer 5 minuten voor het einde de spinazie en de Parmigiano Reggiano
        erdoor, tot de spinazie geslonken is.
  \item Haal de pan van het vuur, roer de boter erdoor en laat de risotto
        heel even rusten. Serveer direct met extra Parmezaanse kaas erover.
\end{steps}
\end{recipe}
